% `tex.web` §1197 normally turns a braced subformula into an Ord noad whose
% nucleus is the list inside the braces. It has one special splice for an
% accent that is the only noad in the group:
%
%     else if type(p)=accent_noad then if saved(0)=nucleus(tail) then
%       if type(tail)=ord_noad then replace the tail of the list by p
%
% The distinction becomes visible when scripts follow the closing brace.
% They attach directly to the accent noad, and §738 consequently swaps the
% accented character and both scripts into the box below the accent. Keeping
% the otherwise invisible outer Ord instead leaves the scripts beside a much
% shorter accent box; this was the first 15.7pt page-10 displacement in TRIP.
\input plain
\showboxdepth=8 \showboxbreadth=40 \tracingonline=0

% The live list must show one accent noad carrying both scripts, not an Ord
% whose braced nucleus contains an accent and whose scripts sit outside it.
\setbox0=\hbox{${\hat A}_b^c\showlists$}
\showbox0

% Scripts already inside the braces do not participate in the splice test,
% but the lone accent noad is still preserved whole when the group closes.
\setbox0=\hbox{${\hat A_b^c}\showlists$}
\showbox0

% A braced accent used as a script remains a sub-mlist. When that list is
% cleaned, however, its one translated vbox is already clean and must be
% returned directly instead of hidden inside a new hbox.
\setbox0=\hbox{$x^{\hat A}\showlists$}
\showbox0

% A group containing anything beside the accent is a real subformula and must
% retain its outer ordinary atom.
\setbox0=\hbox{${\hat A+1}_b^c\showlists$}
\showbox0
\end
